\chapter{Cardiac Catheterization}

\section*{Overview}
Cardiac catheterization involves passing a thin flexible tube (catheter) through a blood vessel in the groin or arm and guiding it into the heart to diagnose or treat cardiovascular conditions. The exact approach, setting, and preparation depend on the indication, anatomy, and the patient's health.

\section*{Alternative Names}
Heart cath, Coronary catheterization, Cardiac cath

\section*{Principle}
A catheter is threaded through a blood vessel into the heart under X-ray guidance, permitting measurement of pressures and oxygen content in the chambers, injection of contrast to image the coronary arteries, and passage of interventional tools such as balloons, stents, and valves.
\section*{History}
Werner Forssmann performed the first human cardiac catheterization on himself in 1929; Andre Cournand and Dickinson Richards developed its clinical use, sharing the 1956 Nobel Prize.

\section*{Why It Is Done}
Ordered to evaluate chest pain, assess heart valve function, measure intracardiac pressures, or locate blockages in coronary arteries.

\section*{Is This a Primary or Secondary Test or Procedure}
Primary diagnostic and guiding intervention for coronary and valvular heart disease.

\section*{How It Is Performed}
Under local anesthesia and mild sedation, a catheter is inserted into the femoral or radial artery and threaded to the heart under fluoroscopic guidance.

\section*{Preparation}
Fasting and medication instructions vary with the indication and sedation plan. Tell the team about anticoagulants, antiplatelet medicines, diabetes medicines, allergies, and kidney disease. Do not stop blood-thinning medicines unless the procedural and prescribing teams give a coordinated plan.

\section*{Contraindications and Precautions}
Active local infection over puncture site; severe uncorrected bleeding disorders.

\section*{Risks}
Bleeding, hematoma at puncture site, cardiac arrhythmias, coronary artery dissection, or contrast allergy.

\section*{Aftercare and Recovery}
Lie flat for 4-6 hours if femoral access was used; monitor puncture site for bleeding; drink fluids to flush contrast.

\section*{Outcome and Follow-up}
Provides detailed information about coronary artery blockages and heart chamber pressures.

\section*{Expected Results}
Normal intracardiac pressures; no significant coronary artery blockages.

\section*{Follow-up}
Coronary angioplasty, CABG, or medical management.

\section*{When to Seek Medical Care}
Follow the procedure team's specific instructions. Seek urgent medical assessment for severe or worsening pain, uncontrolled bleeding, fever or signs of infection, trouble breathing, chest pain, fainting, new weakness or confusion, inability to urinate, or any rapidly worsening symptom. The appropriate warning signs depend on the procedure and the patient's medical condition.

\section*{References}
\begin{enumerate}
  \item MedlinePlus. Cardiac Catheterization. U.S. National Library of Medicine; 2025.
  \item Current Medical Diagnosis and Treatment. McGraw-Hill; 2025.
\end{enumerate}

\clearpage
